% GENERATED FILE. Edit canonical lessons, then run npm run generate:books.
% canonical-source-hash: fnv1a64:cf75fc214b4e25e6
% canonical-lessons: MR-C159-alikade, MR-C159-purvi, MR-C159-adhic, MR-C159-bhutkal, MR-C159-ghadne

\chapter{Recently, Formerly, Already, the Past, To Happen}
\label{ch:mr-recently-formerly-already-the-past-to-happen}
\hlchaptermodality{159}

\begin{chapteropening}
\textbf{By the end of this chapter:} \emph{I can say what happened, and when: recently, already, in the old days.}

Complete the last lesson of chapter 159: I can say what happened, and when: recently, already, in the old days.
\end{chapteropening}

\section[\texorpdfstring{\mr{अलीकडे} (alīkaḍe)}{अलीकडे (alīkaḍe)}]{\mr{अलीकडे} (alīkaḍe) — recently}

\label{lesson:MR-C159-alikade}



\begin{quote}
Before the new one: say the Marathi for to obtain, then the Marathi for to save.
\end{quote}

\subsection*{You'll want to know: \mr{अलीकडे}}
\textbf{\mr{अलीकडे}} — \emph{alīkaḍe} — \textquotedblleft{}recently\textquotedblright{}. Marathi's past ends in \textbf{-\mr{ले}}, \textbf{-\mr{लो}} or \textbf{-\mr{ली}}: \textbf{\mr{अलीकडे} \mr{मी} \mr{एक} \mr{पुस्तक} \mr{वाचले}.} — \emph{alīkaḍe mī ek pustak vācle} — \textquotedblleft{}Recently I read a book.\textquotedblright{}

\begin{grammarlens}[title={What you've built}]
One more word for this chapter.
\end{grammarlens}

\subsection*{Your turn}
\begin{itemize}
\raggedright
  \item \emph{Say it:} \emph{alīkaḍe}
  \item \emph{Say it:} \emph{alīkaḍe}, once more
  \item \emph{Say it:} say \emph{alīkaḍe}
  \item \emph{Recall:} say the Marathi for to obtain, then the Marathi for to save, then say \emph{alīkaḍe} again
\end{itemize}

\begin{tcolorbox}[breakable,colback=teal!4,colframe=teal!35!black,title={Before you move on}]
What does \mr{अलीकडे} mean? (\textquotedblleft{}Recently\textquotedblright{}.) Say it once more.
\end{tcolorbox}
\section[\texorpdfstring{\mr{पूर्वी} (pūrvī)}{पूर्वी (pūrvī)}]{\mr{पूर्वी} (pūrvī) — formerly, before}

\label{lesson:MR-C159-purvi}



\begin{quote}
Before the new one: say the Marathi for to save, then the Marathi for recently.
\end{quote}

\subsection*{You'll want to know: \mr{पूर्वी}}
\textbf{\mr{पूर्वी}} — \emph{pūrvī} — \textquotedblleft{}formerly, before\textquotedblright{}. \textbf{\mr{पूर्वी} \mr{इथे} \mr{एक} \mr{दुकान} \mr{होते}.} — \emph{pūrvī ithe ek dukān hote} — \textquotedblleft{}There used to be a shop here.\textquotedblright{} \textbf{\mr{होते}} is the past of \textbf{\mr{आहे}}.

\begin{grammarlens}[title={What you've built}]
One more word for this chapter.
\end{grammarlens}

\subsection*{Your turn}
\begin{itemize}
\raggedright
  \item \emph{Say it:} \emph{pūrvī}
  \item \emph{Say it:} \emph{pūrvī}, once more
  \item \emph{Say it:} say \emph{pūrvī}
  \item \emph{Recall:} say the Marathi for to save, then the Marathi for recently, then say \emph{pūrvī} again
\end{itemize}

\begin{tcolorbox}[breakable,colback=teal!4,colframe=teal!35!black,title={Before you move on}]
What does \mr{पूर्वी} mean? (\textquotedblleft{}Formerly, before\textquotedblright{}.) Say it once more.
\end{tcolorbox}
\section[\texorpdfstring{\mr{आधीच} (ādhīc)}{आधीच (ādhīc)}]{\mr{आधीच} (ādhīc) — already}

\label{lesson:MR-C159-adhic}



\begin{quote}
Before the new one: say the Marathi for recently, then the Marathi for formerly.
\end{quote}

\subsection*{You'll want to know: \mr{आधीच}}
\textbf{\mr{आधीच}} — \emph{ādhīc} — \textquotedblleft{}already\textquotedblright{}. \textbf{\mr{आधी}}, \textquotedblleft{}before\textquotedblright{}, with the emphatic \textbf{\mr{च}}: \textbf{\mr{मी} \mr{आधीच} \mr{जेवलो}.} — \emph{mī ādhīc jevlo} — \textquotedblleft{}I have already eaten.\textquotedblright{} A woman says \textbf{\mr{जेवले}}.

\begin{grammarlens}[title={What you've built}]
One more word for this chapter.
\end{grammarlens}

\subsection*{Your turn}
\begin{itemize}
\raggedright
  \item \emph{Say it:} \emph{ādhīc}
  \item \emph{Say it:} \emph{ādhīc}, once more
  \item \emph{Say it:} say \emph{ādhīc}
  \item \emph{Recall:} say the Marathi for recently, then the Marathi for formerly, then say \emph{ādhīc} again
\end{itemize}

\begin{tcolorbox}[breakable,colback=teal!4,colframe=teal!35!black,title={Before you move on}]
What does \mr{आधीच} mean? (\textquotedblleft{}Already\textquotedblright{}.) Say it once more.
\end{tcolorbox}
\section[\texorpdfstring{\mr{भूतकाळ} (bhūtkāḷ)}{भूतकाळ (bhūtkāḷ)}]{\mr{भूतकाळ} (bhūtkāḷ) — the past}

\label{lesson:MR-C159-bhutkal}



\begin{quote}
Before the new one: say the Marathi for formerly, then the Marathi for already.
\end{quote}

\subsection*{You'll want to know: \mr{भूतकाळ}}
\textbf{\mr{भूतकाळ}} — \emph{bhūtkāḷ} — \textquotedblleft{}the past\textquotedblright{}. It is \textbf{\mr{भूत}}, \textquotedblleft{}what has been\textquotedblright{}, with \textbf{\mr{काळ}}, \textquotedblleft{}time\textquotedblright{}. Grammar books use it for the past tense too.

\begin{grammarlens}[title={What you've built}]
One more word for this chapter.
\end{grammarlens}

\subsection*{Your turn}
\begin{itemize}
\raggedright
  \item \emph{Say it:} \emph{bhūtkāḷ}
  \item \emph{Say it:} \emph{bhūtkāḷ}, once more
  \item \emph{Say it:} say \emph{bhūtkāḷ}
  \item \emph{Recall:} say the Marathi for formerly, then the Marathi for already, then say \emph{bhūtkāḷ} again
\end{itemize}

\begin{tcolorbox}[breakable,colback=teal!4,colframe=teal!35!black,title={Before you move on}]
What does \mr{भूतकाळ} mean? (\textquotedblleft{}The past\textquotedblright{}.) Say it once more.
\end{tcolorbox}
\section[\texorpdfstring{\mr{घडणे} (ghaḍṇe)}{घडणे (ghaḍṇe)}]{\mr{घडणे} (ghaḍṇe) — to happen}

\label{lesson:MR-C159-ghadne}



\begin{quote}
Before the new one: say the Marathi for already, then the Marathi for the past.
\end{quote}

\subsection*{You'll want to know: \mr{घडणे}}
\textbf{\mr{घडणे}} — \emph{ghaḍṇe} — \textquotedblleft{}to happen\textquotedblright{}. \textbf{\mr{काय} \mr{घडले}?} — \emph{kāy ghaḍle} — \textquotedblleft{}What happened?\textquotedblright{} Its past is \textbf{\mr{घडले}}.

\begin{grammarlens}[title={What you've built}]
One more word for this chapter.
\end{grammarlens}

\subsection*{Your turn}
\begin{itemize}
\raggedright
  \item \emph{Say it:} \emph{ghaḍṇe}
  \item \emph{Say it:} \emph{ghaḍṇe}, once more
  \item \emph{Say it:} say \emph{ghaḍṇe}
  \item \emph{Recall:} say the Marathi for already, then the Marathi for the past, then say \emph{ghaḍṇe} again
\end{itemize}

\begin{tcolorbox}[breakable,colback=teal!4,colframe=teal!35!black,title={Before you move on}]
What does \mr{घडणे} mean? (\textquotedblleft{}To happen\textquotedblright{}.) Say it once more.
\end{tcolorbox}
